\begin{thebibliography}{99}
\bibitem[1]{published-1} Scott Aaronson and Andris Ambainis. The need for structure in quantum speedups. Theory of Computing, 10(6):133–166, 2014.
\bibitem[2]{published-2} Scott Aaronson and Andris Ambainis. Forrelation: A problem that optimally separates quantum from classical computing. SIAM Journal on Computing, 47(3):982–1038, 2018.
\bibitem[3]{published-3} Scott Aaronson, Shalev Ben-David, and Robin Kothari. Separations in query complexity using cheat sheets. In Proceedings of the 48th Annual ACM Symposium on Theory of Computing, pages 863–876, 2016.
\bibitem[4]{published-4} Scott Aaronson, Shalev Ben-David, Robin Kothari, Shravas Rao, and Avishay Tal. Degree vs. approximate degree and quantum implications of Huang’s sensitivity theorem. In Proceedings of the 53rd Annual ACM Symposium on Theory of Computing, pages 1330–1342, 2021. arXiv: 2010.12629.
\bibitem[5]{published-5} Scott Aaronson, DeVon Ingram, and William Kretschmer. The acrobatics of BQP. CoRR, abs/2111.10409, 2021. arXiv:2111.10409.
\bibitem[6]{published-6} Leonard M. Adleman, Jonathan Demarrais, and Ming-Deh A. Huang. Quantum computability. SIAM Journal on Computing, 26(5):1524–1540, 1997.
\bibitem[7]{published-7} Andris Ambainis. Polynomial degree vs. quantum query complexity. Journal of Computer and System Sciences, 72(2):220–238, 2006.
\bibitem[8]{published-8} Andris Ambainis, Kaspars Balodis, Aleksandrs Belovs, Troy Lee, Miklos Santha, and Juris Smotrovs. Separations in query complexity based on pointer functions. Journal of the ACM, 64(5):32:1–32:24, 2017.
\bibitem[9]{published-9} Srinivasan Arunachalam, Jop Briët, and Carlos Palazuelos. Quantum query algorithms are completely bounded forms. SIAM Journal on Computing, 48(3):903–925, 2019.
\bibitem[10]{published-10} Nikhil Bansal and Makrand Sinha. k-Forrelation optimally separates quantum and classical query complexity. In Proceedings of the 53rd Annual ACM SIGACT Symposium on Theory of Computing, pages 1303–1316, 2021.
\bibitem[11]{published-11} Howard Barnum, Michael E. Saks, and Mario Szegedy. Quantum query complexity and semi-definite programming. In Proceedings of 18th Annual IEEE Conference on Computational Complexity, pages 179–193, 2003.
\bibitem[12]{published-12} Robert Beals, Harry Buhrman, Richard Cleve, Michele Mosca, and Ronald de Wolf. Quantum lower bounds by polynomials. Journal of the ACM, 48(4):778–797, 2001.
\bibitem[13]{published-13} Shalev Ben-David, Pooya Hatami, and Avishay Tal. Low-sensitivity functions from unam- biguous certificates. In Proceedings of the 8th Innovations in Theoretical Computer Science Conference, volume 67 of LIPIcs, pages 28:1–28:23. Schloss Dagstuhl - Leibniz-Zentrum für Informatik, 2017.
\bibitem[14]{published-14} Sergey Bravyi, David Gosset, Daniel Grier, and Luke Schaeffer. Classical algorithms for forrelation, 2021. arXiv:2102.06963.
\bibitem[15]{published-15} Harry Buhrman and Ronald de Wolf. Complexity measures and decision tree complexity: a survey. Theoretical Computer Science, 288(1):21–43, 2002.
\bibitem[16]{published-16} Benoît Collins and Camille Male. The strong asymptotic freeness of Haar and deterministic matrices. Annales Scientifiques de l’ENS, (4) 47, fascicule 1:147–163, 2014. arXiv:1105.4345.
\bibitem[17]{published-17} Benoît Collins and Ion Nechita. Random matrix techniques in quantum information theory. Journal of Mathematical Physics, 57(1):015215, 2016.
\bibitem[18]{published-18} Andreas Defant, Mieczyslaw Mastylo, and Antonio Pérez. On the Fourier spectrum of functions on Boolean cubes. Mathematische Annalen, 374:653–680, 2018.
\bibitem[19]{published-19} Irit Dinur, Ehud Friedgut, Guy Kindler, and Ryan O’Donnell. On the Fourier tails of bounded functions over the discrete cube. In Proceedings of the 38th Annual ACM Symposium on Theory of Computing, pages 437–446, 2006.
\bibitem[20]{published-20} Sander Gribling and Monique Laurent. Semidefinite programming formulations for the completely bounded norm of a tensor, 2019. arXiv:1901.04921.
\bibitem[21]{published-21} Lov K. Grover. A fast quantum mechanical algorithm for database search. In Proceedings of the 28th Annual ACM Symposium on Theory of Computing, pages 212–219, 1996.
\bibitem[22]{published-22} Uffe Haagerup. An example of a non nuclear C $\ast$ -algebra, which has the metric approximation property. Inventiones Mathematicae, 50:279–293, 1978/79.
\bibitem[23]{published-23} Todd Kemp and Roland Speicher. Strong Haagerup inequalities for free R-diagonal elements. Journal of Functional Analysis, 251(1):141–173, 2007. arXiv:math/0512481.
\bibitem[24]{published-24} Gatis Midrijanis. On randomized and quantum query complexities, 2005. arXiv:quant-ph/ 0501142.
\bibitem[25]{published-25} Ashley Montanaro. Some applications of hypercontractive inequalities in quantum information theory. Journal of Mathematical Physics, 53(12):122206, 2012.
\bibitem[26]{published-26} Alexandru Nica and Roland Speicher. Lectures on the Combinatorics of Free Probability. London Mathematical Society Lecture Note Series. Cambridge University Press, 2006.
\bibitem[27]{published-27} Noam Nisan and Mario Szegedy. On the degree of Boolean functions as real polynomials. Computational Complexity, 4(4):301–313, 1994.
\bibitem[28]{published-28} Ryan O’Donnell. Analysis of Boolean Functions. Cambridge University Press, 2014.
\bibitem[29]{published-29} Ryan O’Donnell, Michael E. Saks, Oded Schramm, and Rocco A. Servedio. Every decision tree has an influential variable. In Proceedings of the 46th Annual IEEE Symposium on Foundations of Computer Science, pages 31–39, 2005.
\bibitem[30]{published-30} Ryan O’Donnell and Yu Zhao. Polynomial bounds for decoupling, with applications. In Proceedings of 31st Conference on Computational Complexity, volume 50 of LIPIcs, pages 24:1–24:18. Schloss Dagstuhl - Leibniz-Zentrum für Informatik, 2016.
\bibitem[31]{published-31} Alexander A. Sherstov, Andrey A. Storozhenko, and Pei Wu. An optimal separation of randomized and quantum query complexity. In Proceedings of the 53rd Annual ACM SIGACT Symposium on Theory of Computing, pages 1289–1302, 2021.
\bibitem[32]{published-32} Peter W. Shor. Polynomial-time algorithms for prime factorization and discrete logarithms on a quantum computer. SIAM Journal on Computing, 26(5):1484–1509, 1997.
\bibitem[33]{published-33} Daniel R. Simon. On the power of quantum computation. SIAM Journal on Computing, 26(5):1474–1483, 1997.
\bibitem[34]{published-34} Avishay Tal. Towards optimal separations between quantum and randomized query complexities. In Proceedings of the 61st IEEE Annual Symposium on Foundations of Computer Science, pages 228–239, 2020.
\bibitem[35]{published-35} Dan Voiculescu. A strengthened asymptotic freeness result for random matrices with applica- tions to free entropy. International Mathematics Research Notices, 1998:41–63, 1998.
\bibitem[36]{published-36} Z. Yin, A. W. Harrow, M. Horodecki, M. Marciniak, and A. Rutkowski. Random and free observables saturate the Tsirelson bound for CHSH inequality. Physical Review A, 95(032101), 2017. arXiv:1512.00223.
\end{thebibliography}
